\providecommand{\todo}[1]{\textcolor{red}{\textbf{[TODO: #1]}}}
\clearpage
\phantomsection
\section*{Chương 3. Bước làm mềm mép Nyquist: khối FIR tùy chọn}
\addcontentsline{toc}{section}{Chương 3. Bước làm mềm mép Nyquist: khối FIR tùy chọn}

\begin{learningobj}
\begin{itemize}[leftmargin=1.5em, itemsep=2pt]
    \item Thấu hiểu bản chất vật lý và toán học của hiện tượng tràn búp phổ qua tần số Nyquist khi nhân tín hiệu thoại với cửa sổ thời gian hữu hạn, giải mã nguồn gốc dao động đổi dấu luân phiên $(-1)^\tau$ phá hủy đỉnh hàm tự tương quan $r(\tau)$.
    \item Nắm vững lịch sử thuật toán: đối chiếu giải pháp ``làm mềm phổ mép Nyquist'' (soft upsampling) trong bài báo Boersma (1993) với mã nguồn thực tế của Praat 7.0.03; làm rõ sự thật phũ phàng rằng Praat hiện hành không thực hiện bước lọc này trong đường tự tương quan thô (\texttt{raw autocorrelation}).
    \item Khảo sát định lượng 4 phương án thiết kế bộ lọc FIR (Bypass, 31 tap, 63 tap, 127 tap) qua thuật toán bình phương tối thiểu \texttt{firls}; vạch trần sai lầm phổ biến khi nhầm lẫn giữa độ phân giải lượng tử hóa hệ số (16 bit $\approx 96\text{ dB}$) và độ suy hao dải chặn bị khống chế bởi công thức Kaiser.
    \item Phân tích cặn kẽ khái niệm trễ nhóm (group delay) $\tau_g = (N-1)/2$ của bộ lọc FIR đối xứng Type I và quy tắc bắt buộc phải dời lưới khung phân tích 15 mẫu khi kích hoạt bộ lọc.
    \item Khảo cứu toàn diện các tham số của khối IP Intel FIR Compiler II theo tài liệu chuẩn \emph{Intel FIR Compiler II User Guide} (UG-01072); nhận diện cạm bẫy thiết lập tốc độ mẫu \texttt{Input Sample Rate} trong hệ thống đọc tệp.
    \item Thiết kế và hiện thực hóa khối FIR 31 tap gấp thời gian (folded FIR) bằng Verilog-2001, tích hợp bộ tiền cộng (pre-adder) 17 bit và 1 bộ nhân DSP duy nhất chạy 16 chu kỳ/mẫu, tích hợp công tắc chuyển mạch \texttt{bypass} thời gian thực.
    \item Xây dựng mô hình song sinh Python bit-đối-bit và kịch bản kiểm thử chu kỳ với hơn 3.000 chu kỳ kích thích ngẫu nhiên, kiểm chứng đáp ứng xung, sóng sin 7,8 kHz và tín hiệu âm thanh thực tế.
\end{itemize}
\end{learningobj}

\subsection*{3.1. Bản chất hiện tượng rò rỉ phổ mép Nyquist và dao động đổi dấu trong tự tương quan}
\addcontentsline{toc}{subsection}{3.1. Bản chất hiện tượng rò rỉ phổ mép Nyquist và dao động đổi dấu trong tự tương quan}

\subsubsection*{Vấn đề thực tế: Một thành phần gần 8 kHz phá hỏng đỉnh tự tương quan như thế nào?}
Trong hệ thống xử lý tín hiệu tiếng nói với tần số lấy mẫu $f_s = 16.000\text{ Hz}$, tần số giới hạn Nyquist là $f_{\text{Nyq}} = f_s / 2 = 8.000\text{ Hz}$. Theo định lý lấy mẫu Nyquist--Shannon, mọi dao động có tần số thực sự nhỏ hơn $8\text{ kHz}$ đều được bảo toàn trọn vẹn và không bị hiện tượng chồng phổ (aliasing).

Tuy nhiên, trong bài toán ước lượng cao độ $F_0$, chúng ta không bao giờ xử lý một tín hiệu tuần hoàn vô hạn. Thuật toán bắt buộc phải cắt tín hiệu thành từng khung ngắn có độ dài hữu hạn thông qua một hàm cửa sổ thời gian $w[n]$ (cửa sổ Hanning 638 mẫu, thời lượng $T_w \approx 40\text{ ms}$). Về mặt toán học, phép nhân hai dãy tín hiệu trong miền thời gian tương đương với phép tích chập (convolution) giữa phổ của tín hiệu $X(f)$ và phổ của hàm cửa sổ $W(f)$ trong miền tần số liên tục:
\begin{equation}
X_w(f) = X(f) * W(f) = \int_{-f_{\text{Nyq}}}^{f_{\text{Nyq}}} X(\xi) \, W(f - \xi) \, d\xi
\end{equation}

Cửa sổ Hanning có độ rộng búp phổ chính (main lobe) tỷ lệ nghịch với độ dài cửa sổ: $\Delta f \approx 4 / T_w = 4 / 0{,}04\text{ s} = 100\text{ Hz}$, kèm theo chuỗi các búp sóng phụ (side lobes) suy giảm với tốc độ $18\text{ dB/octave}$. Giả sử trong giọng nói xuất hiện một thành phần năng lượng cao tần nằm sát mép phổ, chẳng hạn ở tần số $7{,}8\text{ kHz}$ (họa âm bậc cao của âm thanh thanh môn, phụ âm ma sát vô thanh $/s/, /\int/$, hoặc nhiễu môi trường). Do độ rộng của búp phổ cửa sổ, năng lượng tại $7{,}8\text{ kHz}$ sau tích chập sẽ bị loang rộng sang hai phía. Búp phổ phía trên sẽ tràn vượt qua ranh giới Nyquist $8{,}0\text{ kHz}$, phản xạ ngược vào dải tần số $[0, 8\text{ kHz}]$ và gây ra hiện tượng rò rỉ phổ (spectral leakage).

\subsubsection*{Dẫn xuất toán học: Nguồn gốc sóng dao động đổi dấu luân phiên $(-1)^\tau$}
Hậu quả nghiêm trọng nhất của hiện tượng rò rỉ mép Nyquist không nằm ở việc méo dạng biên độ, mà nằm ở tác động phá hủy cấu trúc của hàm tự tương quan $r(\tau)$.

Hãy xét một dao động hình sin đơn tần nằm đúng tại tần số Nyquist: $f = f_{\text{Nyq}} = 8.000\text{ Hz}$. Chu kỳ thời gian của dao động này là $T_0 = 1 / 8000\text{ s} = 125\text{ }\mu\text{s}$, tương ứng với đúng 2 khoảng lấy mẫu rời rạc ($T_0 / T_s = 2\text{ mẫu}$). Khi được lấy mẫu ở các thời điểm rời rạc $n$, giá trị của tín hiệu có dạng:
\begin{equation}
x[n] = A \cos\left(2\pi \frac{f_{\text{Nyq}}}{f_s} n + \phi_0\right) = A \cos(\pi n + \phi_0)
\end{equation}
Nếu chọn pha ban đầu $\phi_0 = 0$, tín hiệu này trở thành một chuỗi xung luân phiên đổi dấu giữa hai giá trị cực đại:
\begin{equation}
x[n] = A \cdot (-1)^n \quad \implies \quad x = [A, -A, A, -A, A, -A, \dots]
\end{equation}

Khi tính toán hàm tự tương quan thô $R[\tau] = \sum_{n=0}^{N-1-\tau} x[n] x[n+\tau]$ cho độ trễ $\tau$:
\begin{equation}
x[n] \cdot x[n+\tau] = \big(A (-1)^n\big) \cdot \big(A (-1)^{n+\tau}\big) = A^2 (-1)^{2n+\tau} = A^2 \big((-1)^2\big)^n (-1)^\tau = A^2 (-1)^\tau
\end{equation}
Bởi vì $(-1)^{2n} = 1$ với mọi số nguyên $n$, tích của hai mẫu cách nhau $\tau$ bước hoàn toàn độc lập với vị trí thời gian $n$, và chỉ phụ thuộc duy nhất vào tính chẵn lẻ của độ trễ $\tau$:
\begin{equation}
R[\tau] = \sum_{n=0}^{N-1-\tau} A^2 (-1)^\tau = (N - \tau) A^2 \cdot (-1)^\tau = 
\begin{cases} 
+(N - \tau) A^2 & \text{khi } \tau \text{ là số chẵn} \\
-(N - \tau) A^2 & \text{khi } \tau \text{ là số lẻ}
\end{cases}
\end{equation}

Biểu thức trên chỉ ra một sự thật vật lý cốt lõi: \textbf{Mọi thành phần tần số nằm sát mép Nyquist đều tạo ra một sóng dao động đổi dấu liên tục dạng $(-1)^\tau$ trong miền độ trễ $\tau$}.

\subsubsection*{Tính tay trên một chuỗi cụ thể}
Để thấy rõ hiện tượng trên bằng mắt, hãy lấy một đoạn tín hiệu ngắn $N = 4$ mẫu đổi dấu tuần hoàn với biên độ $A = 2$: $x = [2, -2, 2, -2]$. Ta tính lần lượt các giá trị tự tương quan $R[\tau]$ từ $\tau = 0$ đến $\tau = 3$:
\begin{itemize}[leftmargin=1.5em, itemsep=2pt]
    \item Với $\tau = 0$: $R[0] = 2^2 + (-2)^2 + 2^2 + (-2)^2 = 4 + 4 + 4 + 4 = \textbf{+16}$.
    \item Với $\tau = 1$: $R[1] = 2(-2) + (-2)(2) + 2(-2) = -4 - 4 - 4 = \textbf{-12}$.
    \item Với $\tau = 2$: $R[2] = 2(2) + (-2)(-2) = 4 + 4 = \textbf{+8}$.
    \item Với $\tau = 3$: $R[3] = 2(-2) = \textbf{-4}$.
\end{itemize}
Dãy kết quả thu được là $R = [+16, -12, +8, -4]$: biên độ dao động xen kẽ dấu dương--âm với bước nhảy giữa hai lag kế tiếp lên tới 28 đơn vị!

Trong tín hiệu tiếng nói thực tế, tín hiệu dây thanh đới tạo ra một vòm tự tương quan trơn nhẵn với đỉnh cực đại nằm tại chu kỳ cơ bản $\tau_0 \in [27, 213]$ (ví dụ $\tau_0 = 100$ mẫu tương ứng $F_0 = 160\text{ Hz}$). Khi sóng đổi dấu $(-1)^\tau$ của thành phần cao tần chồng lấn lên vòm tự tương quan này, nó tạo thành một dải gai nhọn răng cưa (jagged teeth) biến thiên cực nhanh theo từng bước nhảy $\Delta \tau = 1$. 

Khi thuật toán parabol interpolation ở Chương 9 khớp đường cong bậc hai qua ba điểm lân cận $(R[\tau-1], R[\tau], R[\tau+1])$:
\begin{equation}
\delta = \frac{R[\tau-1] - R[\tau+1]}{2(2R[\tau] - R[\tau-1] - R[\tau+1])}
\end{equation}
Nếu $R[\tau]$ rơi vào lag chẵn, nó bị đội lên một lượng $+A^2$; trong khi hai điểm bên cạnh $\tau \pm 1$ là lag lẻ nên bị kéo tụt xuống một lượng $-A^2$. Mẫu số parabol bị phóng đại nhân tạo, làm đỉnh bị bóp nhọn bất thường. Ngược lại, nếu đỉnh thật rơi vào một lag lẻ, giá trị $R[\tau]$ bị triệt tiêu âm, biến đỉnh cực đại thành một điểm uốn hoặc một thung lũng cục bộ, khiến thuật toán bỏ sót hoàn toàn ứng viên cao độ!

\subsubsection*{Đề xuất của Boersma (1993) và sự thật phũ phàng trong mã nguồn Praat hiện hành}
Để khắc phục hiện tượng răng cưa mép Nyquist nói trên, Paul Boersma trong công trình nền tảng năm 1993 (\emph{Accurate short-term analysis of the fundamental frequency and the harmonics-to-noise ratio of a sampled sound}, Proceedings of the Institute of Phonetic Sciences, Amsterdam, Tập 17, Mục 3.1) đã đề xuất quy trình mang tên ``làm mềm phổ'' (soft upsampling / spectral smoothing):
\begin{quote}
``Trước khi phân tích tự tương quan, phổ của tín hiệu âm thanh được làm mềm tuyến tính từ mức $1{,}0$ tại $95\%$ tần số Nyquist ($7{,}6\text{ kHz}$) dốc đều về $0$ tại $100\%$ tần số Nyquist ($8{,}0\text{ kHz}$) thông qua biến đổi Fourier nhanh FFT.''
\end{quote}
Mục tiêu của hàm dốc $5\%$ dải tần này là triệt tiêu hoàn toàn năng lượng ở sát biên $8\text{ kHz}$ trước khi thực hiện tích chập với cửa sổ Hanning, từ đó ngăn chặn tận gốc hiện tượng tràn phổ và dập tắt sóng đổi dấu $(-1)^\tau$.

\textbf{Sự thật phũ phàng (Sự thật cần nói thẳng):} Khi tiến hành rà soát từng dòng mã nguồn chính thức của Praat các phiên bản 6.4.12 và 7.0.03 (tệp \texttt{fon/Sound\_to\_Pitch.cpp}, dòng TODO-PRAAT-LINE), chúng ta phát hiện một sự thật kỹ thuật mang tính quyết định:
\begin{center}
\textbf{Praat 7.0.03 hoàn toàn KHÔNG thực hiện bất kỳ bước upsampling, resample hay lọc trước làm mềm phổ nào trong luồng xử lý tự tương quan thô (\texttt{To Pitch (raw autocorrelation)...})!}
\end{center}
Boersma đã loại bỏ bước tiền xử lý này trong các bản phát hành chính thức của Praat từ nhiều năm qua. Lý do là vì ở khâu xác định đỉnh tinh xảo, Praat sử dụng thuật toán nội suy \texttt{sinc70} kết hợp tìm kiếm cực trị Brent (\texttt{Sound\_to\_Pitch.cpp}, dòng TODO-PRAAT-LINE). Thuật toán nội suy sinc tái tạo lại hàm tương quan liên tục trong miền thời gian thực, vốn đã có khả năng triệt tiêu các sai số rời rạc hóa mà không cần phải can thiệp làm méo dải tần số cao của tín hiệu gốc.

\subsubsection*{Hệ quả đối soát và quyết định thiết kế trên FPGA}
Sự sai khác giữa bài báo lý thuyết năm 1993 và mã nguồn thực tế dẫn đến một hệ quả kỹ thuật trực tiếp:
\begin{itemize}[leftmargin=1.5em, itemsep=2pt]
    \item Nếu hệ thống FPGA mặc định kích hoạt một khối lọc FIR làm mềm mép Nyquist trong khi phần mềm Praat chuẩn trên máy tính không bật, chúng ta đang \textbf{tự tạo ra một nguồn sai lệch nhân tạo} giữa hai hệ thống. Bộ lọc FIR sẽ làm thay đổi biên độ, làm trôi pha và làm lệch nhẹ các đỉnh tự tương quan của FPGA so với Praat, khiến việc đối soát trở nên bất bình đẳng (khác biệt về thuật toán chứ không phải sai số phần cứng).
    \item Tuy nhiên, việc trang bị khối lọc FIR làm mềm mép Nyquist lại mang giá trị nghiên cứu kỹ thuật to lớn: nó cho phép kỹ sư đo đạc định lượng xem việc triệt tiêu dải tần sát Nyquist có thực sự cải thiện độ chính xác $F_0$ trên các đoạn âm ma sát hay không, và đánh giá chi phí phần cứng của một khối lọc số tiêu chuẩn công nghiệp trên FPGA.
\end{itemize}

\textbf{Quyết định thiết kế (Đã chốt):} Hệ thống trên FPGA sẽ hiện thực một khối lọc FIR 31 tap tùy chọn, tích hợp sẵn đường rẽ nhánh \textbf{\texttt{bypass}}. 
\begin{itemize}[leftmargin=1.5em, itemsep=2pt]
    \item Khi chạy đối soát chuẩn với Praat 7.0.03: khối FIR được \textbf{mặc định cấu hình bypass} (\texttt{bypass = 1}). Dữ liệu âm thanh từ FIFO đi thẳng vào bộ nhớ RAM chia khung mà không qua lọc, đảm bảo đầu vào của FPGA và Praat giống nhau từng bit.
    \item Khi cần nghiên cứu thực nghiệm: người dùng bật cờ \texttt{fir\_bypass = 0} thông qua thanh ghi điều khiển \texttt{ctrl} từ HPS. Lúc này khối FIR 31 tap hoạt động, và hệ thống bắt buộc phải dời lưới định vị khung thời gian lùi lại đúng 15 mẫu để bù trừ trễ nhóm.
\end{itemize}

\textit{Tóm lại: Thành phần gần Nyquist sinh ra sóng đổi dấu $(-1)^\tau$ phá hủy đỉnh tự tương quan; Boersma 1993 đề xuất làm mềm phổ nhưng Praat 7.0.03 đã bỏ bước này; vì vậy khối FIR 31 tap trên FPGA được thiết kế kèm công tắc bypass và mặc định tắt khi đối soát.}

\subsection*{3.2. Các phương án, đánh đổi, ta chọn gì, vì sao (có số liệu kiểm chứng)}
\addcontentsline{toc}{subsection}{3.2. Các phương án, đánh đổi, ta chọn gì, vì sao (có số liệu kiểm chứng)}

\subsubsection*{Vấn đề thực tế: Bài toán đánh đổi giữa độ suy hao, độ trễ và tài nguyên phần cứng}
Khi bắt tay vào hiện thực hóa một bộ lọc số FIR (Finite Impulse Response) trên vi mạch FPGA, câu hỏi đầu tiên đặt ra cho người kỹ sư là: \emph{Cần bao nhiêu tap lọc là đủ, và cái giá phải trả về độ trễ xử lý cũng như tài nguyên phần cứng là gì?}

Một bộ lọc lý tưởng theo đề xuất của Boersma (1993) đòi hỏi đáp ứng biên độ hoàn toàn phẳng ($0\text{ dB}$) từ tần số $0\text{ Hz}$ tới $7{,}6\text{ kHz}$ (tức $95\%$ Nyquist), sau đó đổ dốc thẳng đứng về $0$ tại đúng $8{,}0\text{ kHz}$ (mép Nyquist). Tuy nhiên, nguyên lý bất định trong xử lý tín hiệu số chỉ rõ: một đáp ứng dốc đứng trong miền tần số đòi hỏi một đáp ứng xung dài vô hạn trong miền thời gian ($N \to \infty$). Khi cắt ngắn đáp ứng xung thành $N$ điểm hữu hạn, hiện tượng gợn sóng Gibbs và độ rộng dải chuyển tiếp hữu hạn $\Delta f = 8000 - 7600 = 400\text{ Hz}$ sẽ quyết định trực tiếp chất lượng của bộ lọc.

Để có cơ sở định lượng chính xác, chúng tôi sử dụng thuật toán thiết kế bộ lọc bình phương tối thiểu \texttt{scipy.signal.firls} với mục tiêu đáp ứng biên độ bằng $1{,}0$ trên dải $[0, 7{,}6\text{ kHz}]$ và dốc tuyến tính về $0$ tại $8{,}0\text{ kHz}$. Toàn bộ các thông số kỹ thuật được tính toán và kiểm chứng độc lập thông qua kịch bản kiểm thử \texttt{verify\_v3/ch03\_fir\_design.py}:

\begin{table}[htbp]
\centering
\footnotesize
\setlength{\tabcolsep}{3.5pt}
\caption{So sánh định lượng các phương án bộ lọc FIR làm mềm mép Nyquist tại $f_s = 16\text{ kHz}$}
\label{tab:ch03_fir_options}
\begin{tabular}{lcccccc}
\toprule
Phương án & Số tap ($N$) & Trễ $\tau_g$ & Tại $7{,}4\text{ kHz}$ & Tại $7{,}6\text{ kHz}$ & Tại $8{,}0\text{ kHz}$ & Nhận xét kỹ thuật \\
\midrule
(a) Bypass & 0 & 0 mẫu & $0{,}00\text{ dB}$ & $0{,}00\text{ dB}$ & $0{,}00\text{ dB}$ & Khớp Praat hiện hành; 0 tài nguyên \\
(b) 31 tap & 31 & 15 mẫu & $+0{,}19\text{ dB}$ & $-2{,}29\text{ dB}$ & $-9{,}41\text{ dB}$ & Dải thông phẳng $\pm 0{,}25\text{ dB}$ tới $7{,}4\text{ kHz}$ \\
(c) 63 tap & 63 & 31 mẫu & $+0{,}03\text{ dB}$ & $-0{,}38\text{ dB}$ & $-19{,}54\text{ dB}$ & Suy hao dải biên tăng thêm $\approx 10\text{ dB}$ \\
(d) 127 tap & 127 & 63 mẫu & $-0{,}08\text{ dB}$ & $-0{,}25\text{ dB}$ & $-24{,}20\text{ dB}$ & Độ trễ tăng gấp 4 lần; tốn tài nguyên \\
\bottomrule
\end{tabular}
\end{table}

\begin{figure}[htbp]
\centering
% chuong_v3/fig_ch03_freq_response.tex
% Frequency response comparison of FIR filter options (31, 63, 127 taps) and Boersma 1993 soft transition
\begin{tikzpicture}[>=stealth, font=\small]

    % Styles
    \tikzstyle{gridline}  = [dashed, color=gray!30, line width=0.5pt]
    \tikzstyle{axisline}  = [thick, color=black!80]
    \tikzstyle{curve31}   = [very thick, color=blue!80!black]
    \tikzstyle{curve63}   = [thick, dashed, color=teal!80!black]
    \tikzstyle{curve127}  = [thick, dashdotted, color=purple!80!black]
    \tikzstyle{curvetarget}= [thick, dotted, color=red!80!black, line width=1.0pt]

    % Frame Boundary / Background
    \draw[draw=gray!40, fill=white, rounded corners=6pt, thin] (-1.5, -0.8) rectangle (12.2, 5.0);

    % Title
    \node[anchor=north west, font=\footnotesize\bfseries, color=blue!90!black] at (-1.3, 4.8)
        {Đáp ứng tần số biên độ $|H(f)|$ (dB) của các phương án bộ lọc FIR tại $f_s = 16\text{ kHz}$};

    % Coordinates mapping:
    % x: 0 kHz -> 0.0, 8 kHz -> 10.8 cm (x = (f/1000) * 1.35)
    % y: 0 dB -> 3.5 cm, -10 dB -> 2.3 cm, -20 dB -> 1.1 cm, -30 dB -> -0.1 cm (y = 3.5 + dB * 0.12)

    % Horizontal Grid Lines and Y-Ticks
    \foreach \db/\y/\lbl in {0/3.50/{0}, -5/2.90/{-5}, -10/2.30/{-10}, -15/1.70/{-15}, -20/1.10/{-20}, -25/0.50/{-25}, -30/-0.10/{-30}} {
        \draw[gridline] (0.0, \y) -- (10.8, \y);
        \node[anchor=east, font=\scriptsize, color=gray!90!black] at (-0.1, \y) {\lbl};
    }
    \node[anchor=east, font=\scriptsize\bfseries, color=black!80] at (-0.7, 3.5) {(dB)};

    % Vertical Grid Lines and X-Ticks
    \foreach \f/\x/\lbl in {0/0.0/{0}, 1/1.35/{1}, 2/2.70/{2}, 3/4.05/{3}, 4/5.40/{4}, 5/6.75/{5}, 6/8.10/{6}, 7/9.45/{7}, 7.6/10.26/{7{,}6}, 8/10.80/{8}} {
        \draw[gridline] (\x, -0.1) -- (\x, 3.8);
        \node[anchor=north, font=\scriptsize, color=gray!90!black] at (\x, -0.15) {\lbl};
    }
    \node[anchor=north, font=\scriptsize\bfseries, color=black!80] at (10.8, -0.45) {Tần số (kHz)};

    % Transition band highlight: [7.6 kHz, 8.0 kHz]
    \fill[orange!10, opacity=0.7] (10.26, -0.1) rectangle (10.80, 3.8);
    \node[anchor=south, font=\tiny\bfseries, color=orange!80!black] at (10.53, 3.82) {Dải chuyển 400 Hz};

    % 0 dB reference line (solid)
    \draw[thick, color=gray!60] (0.0, 3.5) -- (10.8, 3.5);

    % Target ideal falloff (Boersma 1993: flat to 7.6 kHz, linear to 0 at 8.0 kHz)
    % 0 dB up to 7.6 kHz (x=10.26), then linear slope down
    \draw[curvetarget] (0.0, 3.5) -- (10.26, 3.5) -- (10.80, 0.0);

    % N = 31 curve
    \draw[curve31] 
        (0.00,3.52) -- (0.27,3.51) -- (0.54,3.49) -- (0.81,3.49) -- (1.08,3.50) --
        (1.35,3.52) -- (1.62,3.51) -- (1.89,3.49) -- (2.16,3.48) -- (2.43,3.50) --
        (2.70,3.52) -- (2.97,3.51) -- (3.24,3.49) -- (3.51,3.48) -- (3.78,3.50) --
        (4.05,3.52) -- (4.32,3.52) -- (4.59,3.49) -- (4.86,3.48) -- (5.13,3.49) --
        (5.40,3.52) -- (5.67,3.52) -- (5.94,3.50) -- (6.21,3.48) -- (6.48,3.49) --
        (6.75,3.52) -- (7.02,3.53) -- (7.29,3.50) -- (7.56,3.47) -- (7.83,3.47) --
        (8.10,3.52) -- (8.37,3.54) -- (8.64,3.51) -- (8.91,3.45) -- (9.18,3.45) --
        (9.45,3.53) -- (9.72,3.60) -- (9.99,3.52) -- (10.26,3.23) -- (10.53,2.72) -- (10.80,2.37);

    % N = 63 curve
    \draw[curve63] 
        (0.00,3.50) -- (0.54,3.50) -- (1.08,3.50) -- (1.62,3.50) -- (2.16,3.50) --
        (2.70,3.50) -- (3.24,3.50) -- (3.78,3.50) -- (4.32,3.50) -- (4.86,3.50) --
        (5.40,3.50) -- (5.94,3.50) -- (6.48,3.50) -- (7.02,3.50) -- (7.56,3.50) --
        (8.10,3.50) -- (8.64,3.50) -- (9.18,3.49) -- (9.72,3.50) -- (9.99,3.50) --
        (10.26,3.45) -- (10.53,2.75) -- (10.80,1.16);

    % N = 127 curve
    \draw[curve127] 
        (0.00,3.50) -- (1.08,3.50) -- (2.16,3.50) -- (3.24,3.50) -- (4.32,3.50) --
        (5.40,3.50) -- (6.48,3.50) -- (7.56,3.50) -- (8.64,3.50) -- (9.45,3.50) --
        (9.99,3.49) -- (10.26,3.47) -- (10.53,2.79) -- (10.80,0.60);

    % Distinct annotations at Nyquist (8.0 kHz)
    \node[anchor=west, font=\tiny\bfseries, color=blue!80!black]   at (10.85, 2.37) {31 tap: $-9{,}4\text{ dB}$};
    \node[anchor=west, font=\tiny\bfseries, color=teal!80!black]   at (10.85, 1.16) {63 tap: $-19{,}5\text{ dB}$};
    \node[anchor=west, font=\tiny\bfseries, color=purple!80!black] at (10.85, 0.60) {127 tap: $-24{,}2\text{ dB}$};

    % Legend Box (placed in empty lower-left area with generous vertical and horizontal spacing)
    \draw[draw=gray!50, fill=white, rounded corners=4pt, thin] (0.4, 0.15) rectangle (6.7, 2.50);
    \node[anchor=north west, font=\scriptsize\bfseries, color=black!85] at (0.6, 2.42) {Chú giải các phương án bộ lọc:};
    \draw[curve31] (0.7, 1.82) -- (1.5, 1.82);
    \node[anchor=west, font=\scriptsize] at (1.6, 1.82) {Phương án (b): 31 tap (chọn, $\tau_g = 15$ mẫu)};
    \draw[curve63] (0.7, 1.37) -- (1.5, 1.37);
    \node[anchor=west, font=\scriptsize] at (1.6, 1.37) {Phương án (c): 63 tap ($\tau_g = 31$ mẫu)};
    \draw[curve127] (0.7, 0.92) -- (1.5, 0.92);
    \node[anchor=west, font=\scriptsize] at (1.6, 0.92) {Phương án (d): 127 tap ($\tau_g = 63$ mẫu)};
    \draw[curvetarget] (0.7, 0.47) -- (1.5, 0.47);
    \node[anchor=west, font=\scriptsize] at (1.6, 0.47) {Boersma (1993) (đường đích lý thuyết)};

    % Axis Frame
    \draw[axisline] (0.0, -0.1) -- (10.8, -0.1);
    \draw[axisline] (0.0, -0.1) -- (0.0, 3.8);

\end{tikzpicture}
\caption{Đáp ứng tần số biên độ $|H(f)|$ (dB) của các phương án bộ lọc FIR so với đường cong mục tiêu Boersma (1993) trên dải chuyển tiếp $400\text{ Hz}$ (trích xuất từ \texttt{verify\_v3/ch03\_fir\_design.py}).}
\label{fig:ch03_freq_response}
\end{figure}

\textbf{Đọc hình.} Hình~\ref{fig:ch03_freq_response} trực quan hóa sự đánh đổi giữa bậc lọc $N$ và độ suy hao tại tần số Nyquist. Đường màu xanh đậm biểu diễn phương án (b) 31 tap: trong toàn bộ dải thoại hữu ích từ $0$ đến $7{,}4\text{ kHz}$, độ gợn biên độ được khống chế chặt chẽ trong khoảng $[-0{,}25\text{ dB}, +0{,}25\text{ dB}]$ (tại $0\text{ kHz}$ gợn $+0{,}13\text{ dB}$, tại $7{,}4\text{ kHz}$ gợn $+0{,}19\text{ dB}$), hầu như không làm biến dạng năng lượng của các formant giọng nói. Khi bước vào dải chuyển tiếp $[7{,}6\text{ kHz}, 8{,}0\text{ kHz}]$ (vùng tô màu cam rộng $400\text{ Hz}$), đáp ứng bắt đầu rơi dốc: đạt $-2{,}29\text{ dB}$ tại $7{,}6\text{ kHz}$ và cán mốc suy hao $-9{,}41\text{ dB}$ tại đúng tần số Nyquist $8{,}0\text{ kHz}$. Đường nét đứt màu xanh ngọc (63 tap) và nét chấm gạch màu tím (127 tap) giúp tăng độ suy hao biên lên lần lượt $-19{,}5\text{ dB}$ và $-24{,}2\text{ dB}$, nhưng đòi hỏi chiều dài thanh ghi dịch và độ trễ tăng vọt.

\subsubsection*{Vạch trần cạm bẫy kỹ thuật: Sai lầm ``31 tap, hệ số 16 bit, suy hao > 60 dB''}
Trong một số tài liệu thiết kế trước đây (điển hình là bản thảo của Argon), người ta thường bắt gặp khẳng định:
\begin{quote}
\emph{``Thiết kế bộ lọc FIR 31 tap với hệ số lượng tử 16 bit cho phép đạt độ suy hao dải chặn lớn hơn 60 dB.''}
\end{quote}
Đây là một khẳng định hoàn toàn sai lầm về mặt lý thuyết xử lý tín hiệu số, xuất phát từ sự nhầm lẫn tai hại giữa hai khái niệm hoàn toàn độc lập:
\begin{enumerate}[leftmargin=1.5em, itemsep=3pt]
    \item \textbf{Độ rộng bit lượng tử hóa hệ số ($B = 16\text{ bit}$):} Tham số này chỉ quyết định \emph{sàn nhiễu lượng tử hóa (quantization noise floor)} do sai số làm tròn số học gây ra. Theo quy tắc xấp xỉ $6{,}02\text{ dB/bit}$, lượng tử hóa 16 bit có dấu cho tỷ số tín hiệu trên nhiễu lượng tử (SQNR) lý thuyết là:
    \begin{equation}
    \text{SQNR} \approx 6{,}02 \times 16 + 1{,}76 \approx 98{,}08\text{ dB}
    \end{equation}
    Độ rộng hệ số 16 bit đảm bảo rằng khi chuyển từ số thực dấu phẩy động sang số nguyên Q1.15, đáp ứng tần số thực tế của bộ lọc bám sát đáp ứng lý thuyết với sai số cực đại nhỏ hơn $0{,}001\text{ dB}$ (đã kiểm chứng trong \texttt{verify\_v3/ch03\_fir\_design.log}). Tuy nhiên, độ rộng bit hệ số \textbf{không tự động tạo ra độ suy hao dải chặn}!
    \item \textbf{Độ suy hao dải chặn ($A_s$) và bề rộng dải chuyển tiếp ($\Delta f$):} Độ suy hao dải chặn hoàn toàn bị chi phối bởi số lượng tap $N$ và tỷ lệ dải chuyển tiếp $\Delta f / f_s$. Mối quan hệ toán học kinh điển này được chuẩn hóa qua công thức ước lượng Kaiser:
    \begin{equation}
    N \approx \frac{A_s - 7{,}95}{14{,}36 \cdot \left(\frac{\Delta f}{f_s}\right)}
    \end{equation}
\end{enumerate}

Hãy áp dụng công thức Kaiser vào bài toán làm mềm mép Nyquist với dải chuyển tiếp chỉ rộng $\Delta f = 8000 - 7600 = 400\text{ Hz}$ ở tần số lấy mẫu $f_s = 16.000\text{ Hz}$ ($\Delta f / f_s = 400 / 16000 = 0{,}025$):
Để đạt được độ suy hao dải chặn mục tiêu $A_s = 60\text{ dB}$, số bậc lọc tối thiểu phải là:
\begin{equation}
N \approx \frac{60 - 7{,}95}{14{,}36 \times 0{,}025} = \frac{52{,}05}{0{,}359} \approx \textbf{144{,}99 tap} \quad \implies \quad N \approx 147\text{ tap}
\end{equation}
(Hàm tối ưu hóa \texttt{scipy.signal.kaiserord(60, 400/8000)} trả về chính xác $N = 147\text{ tap}$ với hệ số cửa sổ $\beta = 5{,}653$).

Như vậy, muốn đạt độ suy hao $60\text{ dB}$ trong khoảng chuyển tiếp hẹp $400\text{ Hz}$, ta cần ít nhất 147 tap! Với chiều dài vỏn vẹn $N = 31\text{ tap}$, việc đòi hỏi độ suy hao $60\text{ dB}$ trong dải chuyển tiếp $400\text{ Hz}$ là điều bất khả thi về mặt vật lý. Bộ lọc 31 tap chỉ có thể tạo ra độ suy hao $-9{,}41\text{ dB}$ tại tần số $8\text{ kHz}$.

\subsubsection*{Trễ nhóm (Group Delay) và quy tắc dịch chuyển lưới khung phân tích}
Một đặc tính cốt lõi của bộ lọc FIR đối xứng Type I ($N = 31$ là số lẻ, $h[n] = h[N - 1 - n]$) là có pha tuyến tính hoàn hảo. Độ trễ nhóm (group delay) là đạo hàm của góc pha theo tần số:
\begin{equation}
\tau_g = -\frac{d\theta(\omega)}{d\omega} = \frac{N - 1}{2}\text{ mẫu}
\end{equation}
Với bộ lọc $N = 31$ tap, độ trễ nhóm là một hằng số tuyệt đối trên mọi tần số:
\begin{equation}
\tau_g = \frac{31 - 1}{2} = \textbf{15 mẫu} \quad \left(= 15 \times \frac{1}{16000}\text{ s} = 0{,}9375\text{ ms}\right)
\end{equation}

\textbf{Cạm bẫy dời lưới khung:} Giả sử khung phân tích thứ $k$ của thuật toán bắt đầu từ mẫu âm thanh $s_{\text{start}}[k]$. Nếu tín hiệu đi qua bộ lọc FIR 31 tap, mỗi mẫu âm thanh đi ra khỏi bộ lọc đã bị trễ thời gian đúng 15 chu kỳ lấy mẫu.
Nếu hệ thống FPGA giữ nguyên chỉ số mẫu bắt đầu $s_{\text{start}}[k]$ mà nạp vào bộ nhớ RAM chia khung, toàn bộ 638 mẫu trong cửa sổ phân tích thực chất là các mẫu của đoạn âm thanh xảy ra trước đó 15 mẫu ($0{,}9375\text{ ms}$) trong tín hiệu gốc!
Sự lệch pha thời gian này sẽ làm dịch chuyển vị trí các đỉnh tự tương quan, làm sai lệch giá trị $F_0$ nội suy và gây ra hiện tượng không khớp ứng viên với kết quả Praat chạy trên tệp WAV gốc.

\textbf{Quy tắc bù trễ bắt buộc:} Khi bật bộ lọc FIR 31 tap (\texttt{bypass = 0}), bộ điều khiển phân khung trên HPS hoặc FPGA bắt buộc phải \textbf{dời lưới khung lùi lại đúng 15 mẫu}:
\begin{equation}
s_{\text{start, FIR}}[k] = s_{\text{start}}[k] - 15
\end{equation}
Nhờ việc định vị khung sớm hơn 15 mẫu, cửa sổ 638 mẫu sau khi đi qua trễ nhóm của bộ lọc FIR sẽ rơi vào đúng vị trí thời gian của khung chưa lọc, bảo toàn tính tương đồng tuyệt đối về mặt vật lý.

\subsubsection*{Lựa chọn thiết kế và các lý do quyết định}
Dựa trên phân tích định lượng ở Bảng~\ref{tab:ch03_fir_options}, chúng tôi lựa chọn hiện thực hóa \textbf{Phương án (b) 31 tap kèm mạch rẽ nhánh bypass} với các lý do kỹ thuật sau:
\begin{enumerate}[leftmargin=1.5em, itemsep=3pt]
    \item \textbf{Độ suy hao vừa đủ và dải thông bảo toàn:} Độ suy hao $-9{,}41\text{ dB}$ tại mép $8\text{ kHz}$ giúp giảm biên độ năng lượng cao tần xuống còn $\approx 33{,}8\%$ (tương đương giảm công suất dao động đổi dấu $(-1)^\tau$ xuống còn $11{,}5\%$), đủ để làm dịu hiện tượng răng cưa mép Nyquist mà vẫn giữ dải thông thoại $[0, 7{,}4\text{ kHz}]$ phẳng trong phạm vi $\pm 0{,}25\text{ dB}$.
    \item \textbf{Độ trễ nhóm nhỏ và dễ bù trừ:} Trễ nhóm 15 mẫu chỉ chiếm $9{,}375\%$ của bước nhảy khung (hop 160 mẫu), hoàn toàn nằm gọn trong dung lượng đệm của bộ nhớ RAM vòng 1024 từ ở Chương 4 mà không đòi hỏi thêm tài nguyên FIFO phụ. Trong khi đó, phương án 127 tap có trễ nhóm lên tới 63 mẫu ($\approx 40\%$ bước nhảy hop), làm tăng đáng kể thời gian đáp ứng thời gian thực.
    \item \textbf{Tối ưu hóa tài nguyên phần cứng:} Nhờ tính đối xứng chẵn, 31 tap lọc được gộp thành 15 cặp cộng trước và 1 hệ số trung tâm. Trong mục 3.4, chúng tôi sẽ chứng minh cấu trúc này chỉ tiêu tốn đúng **1 khối DSP** và **0 khối RAM M10K** trên chip Cyclone V khi áp dụng kỹ thuật gấp thời gian 16 chu kỳ/mẫu.
    \item \textbf{Tương thích kép nhờ mạch bypass:} Việc tích hợp công tắc bypass cho phép đáp ứng cả hai mục tiêu: mặc định bypass để đối soát chuẩn mực từng bit với Praat 7.0.03, và bật lọc khi cần đánh giá ảnh hưởng của hiện tượng mép Nyquist trên thực nghiệm.
\end{enumerate}

\textit{Tóm lại: Bộ lọc 31 tap tạo suy hao $-9{,}4\text{ dB}$ tại Nyquist với trễ nhóm 15 mẫu; muốn suy hao $60\text{ dB}$ trong dải chuyển $400\text{ Hz}$ cần tới 147 tap theo công thức Kaiser; khi bật lọc bắt buộc phải dời lưới khung 15 mẫu.}

\subsection*{3.3. Đường (A): Khối IP FIR Compiler II của Intel}
\addcontentsline{toc}{subsection}{3.3. Đường (A): Khối IP FIR Compiler II của Intel}

\subsubsection*{Vấn đề thực tế: Sự phức tạp của các tham số IP chuẩn công nghiệp}
Khi xây dựng một hệ thống xử lý tín hiệu số trên FPGA Intel, kỹ sư thường ưu tiên sử dụng khối IP có sẵn trong thư viện Intel Quartus Prime mang tên \textbf{FIR Compiler II} (được đặc tả chi tiết trong tài liệu kỹ thuật \emph{Intel FIR Compiler II User Guide}, Mã tài liệu: UG-01072, ấn bản 683111). Khối IP này cung cấp một giao diện đồ họa cấu hình (Parameter Editor) với hàng chục tùy chọn kiến trúc tinh vi.

Tuy nhiên, khối IP là một hộp đen phần cứng: nếu kỹ sư tích chọn các tham số mà không hiểu rõ nguyên lý toán học và cơ chế phân bổ tài nguyên silicon bên dưới, hệ thống sẽ gặp các sự cố nghiêm trọng:
\begin{itemize}[leftmargin=1.5em, itemsep=2pt]
    \item Nếu chọn sai tham số tốc độ lấy mẫu đầu vào (\texttt{Input Sample Rate}), IP sẽ tự động tối ưu hóa phần cứng dựa trên giả định sai, dẫn đến việc hàng đợi bên trong bị tràn và làm rơi rụng các mẫu âm thanh khi chạy ở chế độ tệp.
    \item Nếu chọn sai chế độ làm tròn bit thấp (\texttt{Output LSB Rounding}), tích lũy sai số cắt xén (truncation drift) sẽ tạo ra một điện áp trôi DC nhân tạo, làm biến dạng phép trừ trung bình ở Chương 5.
    \item Nếu không nắm vững thông số \texttt{readyLatency} do IP sinh ra, sự bất đồng tín hiệu điều khiển trên bus Avalon-ST sẽ gây mất mát mẫu dữ liệu như đã phân tích ở Chương 2.
\end{itemize}

Dưới đây, chúng tôi giải mã chi tiết từng tham số cấu hình của khối IP FIR Compiler II cho bài toán lọc làm mềm mép Nyquist 31 tap.

\subsubsection*{Giải thích cặn kẽ từng tham số trong Parameter Editor}

\paragraph{1. Filter Type = Single Rate (Tỷ lệ lấy mẫu đơn).}
Tham số này xác định mối quan hệ giữa tốc độ lấy mẫu ở đầu vào và đầu ra của bộ lọc (\emph{User Guide}, Mục 3.1). Chúng tôi cấu hình \texttt{Single Rate} vì hệ thống đang thực hiện lọc phẳng dải thông mà không thay đổi tần số lấy mẫu ($f_s = 16\text{ kHz}$ ở cả hai đầu). 
Nếu chọn \texttt{Interpolation} (nội suy tăng mẫu) hoặc \texttt{Decimation} (hạ mẫu), IP sẽ chèn thêm số 0 hoặc bỏ bớt mẫu, làm thay đổi tỷ lệ thời gian của tín hiệu âm thanh và phá vỡ cấu trúc khung 638 mẫu của Praat.

\paragraph{2. Symmetry = Symmetrical (Cấu trúc hệ số đối xứng chẵn).}
Hệ số bộ lọc thỏa mãn tính chất đối xứng qua tâm: $h[k] = h[30 - k]$ với $k = 0..14$ và $h[15]$ là hệ số trung tâm. Việc kích hoạt tùy chọn \texttt{Symmetrical} cho phép trình biên dịch IP của Intel tận dụng khối phần cứng \textbf{Pre-adder (Bộ tiền cộng 18-bit)} tích hợp sẵn bên trong các khối DSP chuyên dụng của kiến trúc Cyclone V (\emph{Cyclone V Device Handbook}, Tập 1, Chương 5).
Thay vì phải thực hiện 31 phép nhân riêng biệt, phần cứng sẽ cộng hai mẫu đối xứng trước: $x[n-k] + x[n-(30-k)]$, sau đó mới nhân với hệ số chung $h[k]$. Nhờ vậy, số lượng phép nhân số học cần thực hiện giảm đúng một nửa: từ 31 phép nhân xuống còn $15 + 1 = 16$ phép nhân, tiết kiệm $50\%$ tài nguyên bộ nhân trên chip.

\paragraph{3. Clock Rate = 50.0 MHz.}
Tần số xung nhịp đồng hồ chủ cấp cho khối IP, lấy trực tiếp từ bộ dao động trên bo mạch DE10-Standard. Ở tần số 50 MHz, mỗi chu kỳ xung nhịp có thời lượng đúng $T_{\text{clk}} = 20\text{ ns}$.

\paragraph{4. Input Sample Rate = 0.5 MSPS (Cạm bẫy thiết lập tốc độ dòng dữ liệu).}
Đây là tham số kỹ thuật tiềm ẩn nhiều rủi ro nhất khi cấu hình IP trong hệ thống xử lý tệp.
\begin{itemize}[leftmargin=1.5em, itemsep=2pt]
    \item \emph{Cạm bẫy thiết lập $0{,}016\text{ MSPS}$ ($16\text{ kHz}$):} Rất nhiều người lập trình theo quán tính sẽ điền vào ô này giá trị $0{,}016\text{ MSPS}$ vì tín hiệu âm thanh có tần số lấy mẫu là $16\text{ kHz}$. Tuy nhiên, như đã phân tích sâu sắc ở Chương 2, hệ thống ước lượng cao độ của chúng ta đọc dữ liệu từ tệp WAV trong bộ nhớ RAM thông qua bộ đệm FIFO (file-driven processing), hoàn toàn không bị trói buộc bởi nhịp lấy mẫu $16\text{ kHz}$ của micro ngoại vi! Khi bộ nhớ FIFO có sẵn dữ liệu và khối lọc sẵn sàng, các mẫu âm thanh có thể được chuyển giao với tốc độ hàng triệu mẫu/giây. Nếu ta khai báo $0{,}016\text{ MSPS}$, khối IP FIR Compiler II sẽ tính toán tỷ lệ chia sẻ thời gian nội bộ:
    \begin{equation}
    N_{\text{cycles/sample}} = \frac{f_{\text{clk}}}{f_{\text{sample}}} = \frac{50.000.000\text{ Hz}}{16.000\text{ Hz}} = 3125\text{ chu kỳ/mẫu}
    \end{equation}
    Dựa trên con số 3125 chu kỳ này, bộ sinh IP sẽ thiết kế một bộ ghép kênh thời gian cực đoan (time-multiplexed architecture) với tốc độ tiếp nhận cực kỳ chậm chạp. Khi dòng dữ liệu từ FIFO tuôn vào với tốc độ cao, hàng đợi đầu vào của IP sẽ bị tràn lập tức, gây mất mát mẫu âm thanh mà không phát sinh tín hiệu cảnh báo nào!
    \item \emph{Giải pháp đúng đắn:} Thiết lập \texttt{Input Sample Rate} bằng tốc độ tiếp nhận tối đa thực tế của hệ thống, ví dụ chọn $0{,}5\text{ MSPS}$ (tương ứng với $50\text{ MHz} / 0{,}5\text{ MSPS} = 100\text{ chu kỳ/mẫu}$) hoặc $1{,}0\text{ MSPS}$ ($50\text{ chu kỳ/mẫu}$). Cấu hình này ép buộc IP phải bố trí đủ phần cứng để có thể hoàn tất việc tiếp nhận và xử lý một mẫu mới trong vòng 100 chu kỳ xung nhịp, hoàn toàn đáp ứng thông lượng dòng dữ liệu và tương thích với nhịp độ xử lý của khối MAC phía sau ($\approx 616\text{ chu kỳ/mẫu}$).
\end{itemize}

\paragraph{5. Coefficient Width = 16 bit, Coefficient Scaling = Auto/Normalized.}
Độ rộng bit của hệ số lọc được chọn là 16-bit có dấu (signed 2's complement). Tùy chọn \texttt{Auto/Normalized} tự động nhân toàn bộ đáp ứng xung với hệ số tỷ lệ $2^{15} = 32768$ (định dạng Q1.15), đưa hệ số trung tâm $h[15] = 0{,}975000$ lên giá trị nguyên $31949$ (\texttt{0x7CCD}). Việc chuẩn hóa này tận dụng tối đa dải động 16-bit mà không gây tràn số học.

\paragraph{6. Input Data Type = Signed Binary, Width = 16 bit, Fractional Bits = 15.}
Định dạng số học đầu vào khớp chính xác với định dạng mẫu âm thanh PCM 16-bit có dấu đọc từ tệp WAV (khoảng giá trị $[-32768, 32767]$).

\paragraph{7. Output MSB/LSB Rounding và Output Bit Width.}
Trong quá trình tính toán tích chập, mỗi phép nhân giữa mẫu 16-bit và hệ số 16-bit sinh ra tích 32-bit (với bộ tiền cộng 17-bit thì tích là 33-bit). Khi cộng dồn 31 tích (tương ứng mở rộng $\lceil \log_2(31) \rceil = 5\text{ bit}$), giá trị tích lũy nội bộ có thể đạt độ rộng tới 38-bit.
Để ngõ ra đưa vào bộ nhớ RAM chia khung ở Chương 4 giữ nguyên định dạng mẫu âm thanh 16-bit có dấu:
\begin{itemize}[leftmargin=1.5em, itemsep=2pt]
    \item Ta cấu hình \texttt{Output Bit Width} = 16 bit signed.
    \item \texttt{Output LSB Rounding} = \textbf{Convergent Rounding (Làm tròn hội tụ về số chẵn gần nhất)} hoặc \textbf{Nearest (Làm tròn về số nguyên gần nhất)}. Tùy chọn này cộng một nửa trọng số LSB ($2^{14} = 16384$) trước khi dịch phải 15 bit, triệt tiêu hoàn toàn sai số trôi DC so với chế độ \texttt{Truncation} (cắt cụt).
    \item \texttt{Output MSB Rounding} = \textbf{Saturation (Bão hòa)}. Trong trường hợp tín hiệu âm thanh có các đỉnh biên độ cực hạn kết hợp với hiện tượng gợn sóng vượt ngưỡng (overshoot) của bộ lọc khiến giá trị làm tròn vượt quá $32767$, bộ bão hòa sẽ ghim giá trị tại $32767$ (thay vì để tràn bit bù hai biến thành số âm lớn $-32768$, gây méo dạng tiếng nổ cực kỳ nghiêm trọng).
\end{itemize}

\paragraph{8. Back-pressure Support = Yes.}
Kích hoạt giao thức bắt tay hai chiều Avalon-ST thông qua cặp tín hiệu \texttt{ready} và \texttt{valid}. Điều này đảm bảo khi hạ nguồn (khối RAM chia khung hoặc khối tự tương quan) bị bận, tín hiệu phản áp \texttt{ready} hạ xuống 0 sẽ lập tức tạm dừng luồng dữ liệu của bộ lọc FIR một cách an toàn mà không làm mất mẫu.

\subsubsection*{Những gì khối IP sinh ra và giao diện ghép nối phần cứng}
Sau khi hoàn tất quá trình sinh mã trong Platform Designer hoặc IP Catalog, khối IP FIR Compiler II tạo ra một mô-đun phần cứng với các giao tiếp chuẩn sau:
\begin{itemize}[leftmargin=1.5em, itemsep=2pt]
    \item \textbf{Cổng điều khiển:} xung nhịp \texttt{clk} (50 MHz) và tín hiệu xóa \texttt{reset\_n}.
    \item \textbf{Giao diện nhận (Avalon-ST Sink):} bus dữ liệu \texttt{ast\_sink\_data} (16 bit), cờ \texttt{ast\_sink\_valid}, cờ sẵn sàng \texttt{ast\_sink\_ready}, và cờ báo lỗi \texttt{ast\_sink\_error}.
    \item \textbf{Giao diện phát (Avalon-ST Source):} bus dữ liệu \texttt{ast\_source\_data} (16 bit), cờ \texttt{ast\_source\_valid}, tín hiệu nhận \texttt{ast\_source\_ready}, và cờ báo lỗi \texttt{ast\_source\_error}.
    \item \textbf{Độ trễ xử lý (Latency):} Báo cáo sinh mã (IP Generation Report) của Intel FIR Compiler II sẽ công bố chính xác số chu kỳ xung nhịp từ khi mẫu đầu tiên được nạp vào đến khi mẫu đầu ra xuất hiện. Thông số độ trễ này phụ thuộc vào mức độ đường ống hóa nội bộ mà IP tự động phân bổ, và \textbf{phải kiểm trong mô phỏng IP hoặc báo cáo sinh mã của Quartus}.
    \item \textbf{Tài nguyên phần cứng chiếm dụng:} Số lượng khối logic ALM, số khối nhân DSP và bộ nhớ M10K phụ thuộc vào phiên bản Quartus Prime và thuật toán Fitter tối ưu hóa, do đó con số cụ thể là \textbf{(phải đo)}.
\end{itemize}

Dưới đây là đoạn mã Verilog giao diện chuẩn (instantiation template) để tích hợp khối IP FIR Compiler II vào hệ thống:

\begin{lstlisting}[language=Verilog, caption={Mã khởi tạo giao diện khối IP Intel FIR Compiler II trong hệ thống}]
fir_compiler_ii_wrapper u_fir_ip (
    .clk              (clk),
    .reset_n          (rst_n),
    // Avalon-ST Sink Interface (Dữ liệu vào từ FIFO Unpacker)
    .ast_sink_data    (fir_in_data),   // signed [15:0]
    .ast_sink_valid   (fir_in_valid),  // 1-bit
    .ast_sink_ready   (fir_in_ready),  // 1-bit phản áp về FIFO
    .ast_sink_error   (2'b00),
    // Avalon-ST Source Interface (Dữ liệu ra tới RAM khung)
    .ast_source_data  (fir_out_data),  // signed [15:0]
    .ast_source_valid (fir_out_valid), // 1-bit
    .ast_source_ready (fir_out_ready), // 1-bit phản áp từ hạ nguồn
    .ast_source_error ()
);
\end{lstlisting}

\textit{Tóm lại: Khối IP Intel FIR Compiler II tối ưu hóa tài nguyên nhờ đối xứng Pre-adder; cần đặc biệt cảnh giác đặt Input Sample Rate theo thông lượng thực của tệp (0,5 MSPS) thay vì tốc độ danh định 16 kHz; tài nguyên và độ trễ chính xác của IP phải kiểm trong mô phỏng Quartus.}

\subsection*{3.4. Đường (B): Khối HDL tự viết: Bộ lọc FIR đối xứng gấp thời gian}
\addcontentsline{toc}{subsection}{3.4. Đường (B): Khối HDL tự viết: Bộ lọc FIR đối xứng gấp thời gian}

\subsubsection*{Bản chất kiến trúc: Gấp thời gian (Folding) và Bộ tiền cộng (Pre-adder)}
Trong khi khối IP FIR Compiler II đóng vai trò như một giải pháp hộp đen tiện lợi nhưng khó tùy biến sâu, việc tự tay thiết kế bộ lọc FIR bằng ngôn ngữ phần cứng Verilog-2001 đem lại sự minh bạch hoàn toàn về cấu trúc mạch điện, giúp người kỹ sư làm chủ từng bit tính toán và tối ưu hóa diện tích chip tới mức triệt để.

Đối với một bộ lọc FIR thông thường có bậc $N = 31$, công thức tích chập trực tiếp là:
\begin{equation}
y[n] = \sum_{i=0}^{30} h[i] \cdot x[n - i]
\end{equation}
Nếu hiện thực theo kiến trúc song song trực tiếp (Fully Parallel Direct Form), mạch điện đòi hỏi tới $31$ bộ nhân 16-bit và một cây cộng 31 nhánh. Kiến trúc này có thể xử lý 1 mẫu mỗi chu kỳ (thông lượng 50 MSPS ở 50 MHz), nhưng tiêu tốn tới 16 khối DSP phần cứng trên chip Cyclone V.

Tuy nhiên, trong hệ thống ước lượng cao độ $F_0$, tốc độ dòng âm thanh thực tế chỉ là $16\text{ kHz}$ (tương đương với $3125$ chu kỳ xung nhịp 50 MHz cho mỗi mẫu), và ngay cả ở chế độ xử lý tệp nhanh nhất, khối tự tương quan MAC ở Chương 7 cũng cần tới $\approx 616\text{ chu kỳ}$ để xử lý một mẫu âm thanh. Việc sử dụng tới 16 bộ nhân DSP để tính một bộ lọc chỉ cần chạy 1 lần trong hàng trăm chu kỳ là một sự lãng phí tài nguyên silicon cực kỳ lớn!

Để giải quyết mâu thuẫn này, chúng tôi áp dụng kiến trúc **gấp thời gian (Time-Multiplexed / Folded FIR Architecture)** kết hợp với **bộ tiền cộng đối xứng (Symmetric Pre-adder)**:
\begin{enumerate}[leftmargin=1.5em, itemsep=2pt]
    \item Do đáp ứng xung đối xứng chẵn qua tâm $h[i] = h[30 - i]$, ta nhóm các tích có chung hệ số lại với nhau:
    \begin{equation}
    y[n] = \sum_{k=0}^{14} h[k] \cdot \big(x[n - k] + x[n - (30 - k)]\big) + h[15] \cdot x[n - 15]
    \end{equation}
    Tổng cộng chỉ còn đúng 16 số hạng cần nhân.
    \item Phép cộng hai mẫu âm thanh $x[n-k] + x[n-(30-k)]$ được thực hiện trước bằng một bộ cộng tổ hợp 17-bit mang tên \textbf{Pre-adder}. Vì mỗi mẫu âm thanh là số có dấu 16-bit (khoảng $[-32768, 32767]$), tổng của hai mẫu cần đúng 17-bit có dấu (khoảng $[-65536, 65534]$) để không bao giờ bị tràn số.
    \item Thay vì dùng 16 bộ nhân phần cứng riêng biệt, ta sử dụng **đúng 1 bộ nhân DSP duy nhất** và cho nó chạy tuần tự qua 16 bước xung nhịp ($k = 0..15$). Tại mỗi bước xung nhịp, bộ nhân tính một tích 33-bit và cộng dồn kết quả vào một thanh ghi tích lũy 40-bit.
\end{enumerate}

Sau đúng 16 chu kỳ tính toán, kết quả tích lũy hoàn chỉnh được đưa qua khâu làm tròn bán phần (+16384) và bão hòa 16-bit để xuất ra ngõ ra. Toàn bộ kiến trúc phần cứng này được minh họa chi tiết trên Hình~\ref{fig:ch03_folded_arch}.

\begin{figure}[htbp]
\centering
% chuong_v3/fig_ch03_folded_arch.tex
% Architecture of the 31-tap folded symmetric FIR filter with 17-bit pre-adder and 1 DSP block
\begin{tikzpicture}[scale=0.82, transform shape, >=stealth, font=\small]

    % Style definitions
    \tikzstyle{block}     = [draw=blue!80!black, fill=blue!6, thick, rounded corners=3pt, 
                             inner sep=5pt, align=center, font=\scriptsize\bfseries]
    \tikzstyle{regbox}    = [draw=teal!80!black, fill=teal!8, thick, rounded corners=3pt, 
                             inner sep=5pt, align=center, font=\scriptsize\bfseries]
    \tikzstyle{dspbox}    = [draw=red!80!black, fill=red!6, thick, rounded corners=4pt, 
                             inner sep=6pt, align=center, font=\scriptsize\bfseries]
    \tikzstyle{rombox}    = [draw=purple!80!black, fill=purple!6, thick, rounded corners=3pt, 
                             inner sep=5pt, align=center, font=\scriptsize\bfseries]
    \tikzstyle{muxbox}    = [draw=orange!80!black, fill=orange!12, thick, 
                             inner sep=3pt, align=center, font=\scriptsize\bfseries]
    \tikzstyle{fsmbox}    = [draw=gray!80!black, fill=gray!8, thick, rounded corners=4pt, 
                             inner sep=5pt, align=center, font=\scriptsize\bfseries]
    \tikzstyle{dataline}  = [->, very thick, color=teal!80!black]
    \tikzstyle{ctrlline}  = [->, thick, dashed, color=purple!80!black]
    \tikzstyle{portlabel} = [font=\scriptsize\ttfamily]

    % Outer Boundary
    \draw[draw=blue!60!black, fill=blue!1, thick, rounded corners=6pt, dashed] 
        (-4.8, -4.7) rectangle (9.2, 4.0);
    \node[anchor=north west, font=\footnotesize\bfseries, color=blue!90!black] 
        at (-4.6, 3.85) {Khối lọc FIR 31 tap gấp thời gian (\texttt{ch03\_fir\_filter.v})};

    % Upstream Avalon-ST Sink Ports (Left)
    \node[portlabel, anchor=east] (sink_data)  at (-5.0,  1.6) {sink\_data[15:0]};
    \node[portlabel, anchor=east] (sink_valid) at (-5.0, -1.6) {sink\_valid};
    \node[portlabel, anchor=east] (sink_ready) at (-5.0, -2.4) {sink\_ready};
    \node[portlabel, anchor=east] (bypass_in)  at (-5.0, -3.3) {bypass};

    % Downstream Avalon-ST Source Ports (Right)
    \node[portlabel, anchor=west] (src_data)   at (9.4,  1.6) {source\_data[15:0]};
    \node[portlabel, anchor=west] (src_valid)  at (9.4, -1.6) {source\_valid};
    \node[portlabel, anchor=west] (src_ready)  at (9.4, -2.4) {source\_ready};

    % 1. Shift Register Block (31 stages x 16 bit)
    \node[regbox, minimum width=2.2cm, minimum height=1.3cm] (shift_reg) at (-3.1, 1.6)
        {Mảng thanh ghi dịch\\31 tầng $\times$ 16 bit\\(\texttt{shift\_reg[0..30]})};

    % 2. Pre-adder (17-bit signed)
    \node[block, minimum width=1.8cm, minimum height=1.3cm] (preadder) at (-0.3, 1.6)
        {Bộ tiền cộng\\Pre-adder 17 bit\\$x[k] + x[30-k]$};

    % 3. Coefficient ROM (16 x 16 bit)
    \node[rombox, minimum width=1.8cm, minimum height=1.1cm] (rom) at (-0.3, -0.6)
        {ROM hệ số Q1.15\\16 từ $\times$ 16 bit\\(\texttt{h[0..15]})};

    % 4. Multiplier (17 x 16 -> 33 bit inside DSP)
    \node[dspbox, minimum width=2.2cm, minimum height=1.3cm] (dsp_mult) at (2.3, 1.6)
        {Bộ nhân DSP\\17-bit $\times$ 16-bit\\Tích $p$ (33 bit)};

    % 5. Accumulator (40-bit)
    \node[regbox, minimum width=2.0cm, minimum height=1.1cm] (acc) at (2.3, -0.6)
        {Bộ tích lũy\\Thanh ghi 40 bit\\(\texttt{acc})};

    % 6. Rounding & Saturation
    \node[block, minimum width=2.2cm, minimum height=1.3cm] (round_sat) at (5.8, 1.6)
        {Làm tròn \& Bão hòa\\$(acc + 16384) \gg 15$\\Ghim $[-32768, 32767]$};

    % 7. Bypass Multiplexer
    \node[muxbox, minimum width=0.8cm, minimum height=1.5cm] (bypass_mux) at (7.8, 1.6)
        {MUX\\2:1};

    % 8. FSM Controller & Step Counter
    \node[fsmbox, minimum width=6.2cm, minimum height=1.1cm] (fsm) at (1.8, -3.2)
        {Bộ điều khiển FSM \& Bộ đếm bước 4-bit\\Trạng thái: \texttt{S\_IDLE}, \texttt{S\_CALC}, \texttt{S\_WAIT\_OUT}\\Đếm \texttt{step} = $0..15$ (16 chu kỳ/mẫu)};

    % -------------------------------------------------------------
    % Data Connections
    % -------------------------------------------------------------
    % Input data to Shift Register
    \draw[dataline] (-4.8, 1.6) -- (shift_reg.west);

    % Bypass direct line (taps at x = -4.3, routes above all blocks)
    \fill (-4.3, 1.6) circle (2pt);
    \draw[dataline, color=orange!85!black] 
        (-4.3, 1.6) -- (-4.3, 3.2) -- (7.8, 3.2) -- (bypass_mux.north);
    \node[below, font=\scriptsize\bfseries, color=orange!90!black] 
        at (5.4, 3.1) {Đường Bypass trực tiếp (bypass = 1)};
    \node[font=\scriptsize\bfseries, color=orange!90!black, anchor=south east] 
        at (7.7, 2.4) {1};

    % Shift Register to Pre-adder
    \draw[dataline] (shift_reg.east) -- (preadder.west);
    \node[above=4pt, font=\scriptsize] at (-1.5, 1.6) {Cặp mẫu};

    % Pre-adder to DSP Multiplier (input A)
    \draw[dataline] (preadder.east) -- ([yshift=0.2cm]dsp_mult.west);
    \node[above=4pt, font=\scriptsize] at (0.9, 1.8) {17 bit};

    % ROM to DSP Multiplier (input B)
    \draw[dataline] (rom.east) -- ++(0.3, 0) -- ++(0, 1.9) -- ([yshift=-0.2cm]dsp_mult.west);
    \node[anchor=west, font=\scriptsize] at (1.0, 0.4) {$h[step]$ (16b)};

    % DSP Multiplier to Accumulator (p input)
    \draw[dataline] ([xshift=-0.4cm]dsp_mult.south) -- ([xshift=-0.4cm]acc.north);
    \node[left, font=\scriptsize] at (1.9, 0.5) {Tích $p$};

    % Accumulator feedback to Multiplier/Adder
    \draw[dataline] ([xshift=0.4cm]acc.north) -- ([xshift=0.4cm]dsp_mult.south);
    \node[right, font=\scriptsize] at (2.7, 0.5) {Cộng dồn};

    % Accumulator to Rounding & Saturation
    \draw[dataline] (acc.east) -- (4.0, -0.6) -- (4.0, 1.6) -- (round_sat.west);
    \node[right, font=\scriptsize] at (4.05, 0.5) {acc (40b)};

    % Rounding & Saturation to Bypass MUX (input 0)
    \draw[dataline] (round_sat.east) -- (bypass_mux.west);
    \node[above=4pt, font=\scriptsize] at (7.3, 1.6) {Lọc (0)};

    % MUX output to source_data port
    \draw[dataline] (bypass_mux.east) -- (9.2, 1.6);

    % -------------------------------------------------------------
    % Control Connections (Dashed purple)
    % -------------------------------------------------------------
    % sink_valid input to FSM
    \draw[ctrlline] (-4.8, -1.6) -- (-1.8, -1.6) -- (-1.8, -2.65);
    \node[above, font=\scriptsize] at (-3.2, -1.6) {\texttt{sink\_valid}};

    % sink_ready output from FSM
    \draw[ctrlline] (-1.3, -2.65) -- (-1.3, -2.4) -- (-4.8, -2.4);
    \node[above, font=\scriptsize] at (-3.2, -2.4) {\texttt{sink\_ready}};

    % bypass input to FSM and MUX
    \draw[ctrlline] (-4.8, -3.3) -- (fsm.west);
    \fill (-3.6, -3.3) circle (1.8pt);
    \draw[ctrlline, color=orange!80!black] 
        (-3.6, -3.3) -- (-3.6, -4.4) -- (7.8, -4.4) -- (bypass_mux.south);
    \node[above=2pt, font=\scriptsize, color=orange!90!black] 
        at (2.0, -4.4) {\texttt{bypass} (chọn kênh MUX 2:1)};

    % source_ready input to FSM
    \draw[ctrlline] (9.2, -2.4) -- (4.5, -2.4) -- (4.5, -2.65);
    \node[above, font=\scriptsize] at (6.8, -2.4) {\texttt{source\_ready}};

    % source_valid output from FSM
    \draw[ctrlline] (5.1, -2.65) -- (5.1, -1.6) -- (9.2, -1.6);
    \node[above, font=\scriptsize] at (6.8, -1.6) {\texttt{source\_valid}};

    % step[3:0] counter to ROM
    \draw[ctrlline] (-0.3, -2.65) -- (rom.south);
    \node[right, font=\scriptsize] at (-0.3, -1.6) {\texttt{step[3:0]}};

    % latch_out control from FSM to Round_Sat
    \draw[ctrlline] (3.6, -2.65) -- (3.6, -1.4) -- (5.8, -1.4) -- (round_sat.south);
    \node[above, font=\scriptsize] at (4.7, -1.4) {\texttt{latch\_out}};

\end{tikzpicture}

\caption{Sơ đồ kiến trúc phần cứng khối lọc FIR 31 tap gấp thời gian với bộ tiền cộng 17-bit và 1 khối DSP (\texttt{rtl/v3/ch03\_fir\_filter.v}).}
\label{fig:ch03_folded_arch}
\end{figure}

\textbf{Đọc hình.} Hình~\ref{fig:ch03_folded_arch} mô tả chi tiết dòng chảy dữ liệu bên trong mô-đun phần cứng:
\begin{itemize}[leftmargin=1.5em, itemsep=2pt]
    \item Phía bên trái là giao tiếp Avalon-ST Sink: khi tín hiệu bắt tay \texttt{s\_handshake} xuất hiện, mẫu mới \texttt{sink\_data} được nạp vào đầu mảng thanh ghi dịch 31 tầng (tạo thành từ 496 flip-flop trong ALM).
    \item Bộ tiền cộng Pre-adder 17-bit rút các cặp mẫu đối xứng $x[k]$ và $x[30-k]$ từ thanh ghi dịch theo chỉ số bước \texttt{step} do bộ đếm FSM điều khiển. Tại bước cuối cùng (\texttt{step} = 15), ngõ vào thứ hai được nối đất về 0 để chỉ nhân hệ số trung tâm $h[15] \cdot x[15]$.
    \item Bảng ROM hệ số 16 từ phát ra giá trị $h[step]$ tương ứng. Bộ nhân DSP thực hiện phép nhân $17\text{ bit} \times 16\text{ bit} \to 33\text{ bit}$, nạp vào bộ cộng tích lũy 40-bit.
    \item Tại bước 15, giá trị tổng được làm tròn và bão hòa rồi nạp vào bộ dồn kênh MUX 2:1. Tín hiệu điều khiển \texttt{bypass} cho phép dữ liệu từ ngõ vào đi thẳng ra ngõ ra qua 1 tầng chốt đồng bộ khi ở chế độ chuyển tiếp.
\end{itemize}

\subsubsection*{Cạm bẫy kỹ thuật số học: Độ chính xác trung gian và quy tắc làm tròn}
Để phần cứng tự viết khớp từng bit với khối IP công nghiệp và mô hình toán học, lập trình viên phải giải quyết hai cạm bẫy số học tinh vi:
\begin{itemize}[leftmargin=1.5em, itemsep=3pt]
    \item \textbf{Cạm bẫy độ rộng bộ tích lũy (Accumulator Width):} 
    Tích sau bộ nhân có độ rộng 33 bit có dấu. Khi cộng dồn 16 tích số có độ lớn cực đại (trường hợp tín hiệu xoay chiều biên độ tối đa), giá trị tổng cần thêm $\lceil \log_2(16) \rceil = 4\text{ bit}$ mở rộng dấu, nâng độ rộng tối thiểu lên $33 + 4 = 37\text{ bit}$. Nếu lập trình viên chỉ khai báo thanh ghi 32-bit hoặc 36-bit, bộ tích lũy sẽ bị tràn số bit cao trong các đoạn âm thanh hô lớn hoặc tiếng thở mạnh. Trong thiết kế này, chúng tôi chọn độ rộng bộ tích lũy là **40-bit signed** (\texttt{reg signed [39:0] acc}), đảm bảo một khoảng dự trữ an toàn 3 bit chống tràn trong mọi tình huống cực đoan.
    \item \textbf{Cạm bẫy làm tròn đối xứng (Symmetric Round-to-Nearest):}
    Vì các hệ số lọc được định dạng ở chuẩn Q1.15 (nhân với $2^{15} = 32768$), giá trị tích lũy nội bộ bị phóng đại $2^{15}$ lần. Để đưa dữ liệu trở về chuẩn số nguyên 16-bit ban đầu, ta phải thực hiện phép chia cho $2^{15}$ (tức dịch phải số học 15 bit, ký hiệu \texttt{>>> 15}).
    Nếu thực hiện dịch phải trực tiếp, phép toán tương đương với làm tròn hướng về âm vô cùng (floor), tạo ra sai số thiên lệch âm trung bình $-0{,}5\text{ LSB}$. Để làm tròn về số nguyên gần nhất (Round-to-Nearest), ta bắt buộc phải cộng thêm nửa trọng số của bit LSB bị cắt bỏ, tức cộng thêm hằng số $2^{14} = 16384$:
    \begin{equation}
    \text{acc}_{\text{rounded}} = (\text{acc}_{\text{final}} + 16384) \ggg 15
    \end{equation}
    \item \textbf{Khâu bão hòa bảo vệ (Clamping):}
    Mặc dù tổng các hệ số lọc lý thuyết xấp xỉ bằng 1, đáp ứng tần số thực tế có độ gợn nhẹ $+0{,}19\text{ dB}$ ở vùng dải thông cao (tương đương độ lợi biên độ $1{,}022$). Do đó, một tín hiệu vào có biên độ cực đại $32767$ khi đi qua bộ lọc có thể sinh ra giá trị làm tròn lên tới $33256$ (như đã chứng minh trong bài tập ở mục 3.7). Trong định dạng bù hai 16-bit, giá trị $33256$ sẽ bị cuộn tràn thành $-32280$, làm lật ngược dấu của tín hiệu và tạo ra tiếng nổ lách tách phá hủy toàn bộ phép tính $F_0$.
    Khâu bão hòa được cài đặt chặt chẽ:
    \begin{equation}
    y[n] = 
    \begin{cases}
    +32767 & \text{khi } \text{acc}_{\text{rounded}} > +32767 \\
    -32768 & \text{khi } \text{acc}_{\text{rounded}} < -32768 \\
    \text{acc}_{\text{rounded}}[15:0] & \text{trong các trường hợp còn lại}
    \end{cases}
    \end{equation}
\end{itemize}

\subsubsection*{Mã nguồn Verilog hiện thực hóa hoàn chỉnh}
Dưới đây là toàn bộ mã nguồn phần cứng của khối lọc FIR 31 tap gấp thời gian (\texttt{rtl/v3/}\\ \texttt{ch03\_fir\_filter.v}):

\lstinputlisting[language=Verilog, firstline=22, lastline=195, caption={Mã nguồn Verilog hoàn chỉnh của khối FIR 31 tap gấp thời gian (rtl/v3/ch03\_fir\_filter.v)}]{rtl/v3/ch03_fir_filter.v}

\subsubsection*{Giải thích mã nguồn dòng-từng-dòng}
Để bạn đọc nắm bắt tường tận cơ chế vận hành của mạch, chúng tôi thuyết minh chi tiết từng khối mã trong tệp \texttt{rtl/v3/ch03\_fir\_filter.v}:
\begin{itemize}[leftmargin=1.5em, itemsep=2pt]
    \item \textbf{Dòng 22--40:} Khai báo module \texttt{ch03\_fir\_filter} với 3 tham số generic: độ rộng dữ liệu âm thanh \texttt{DATA\_WIDTH = 16}, độ rộng hệ số lọc \texttt{COEFF\_WIDTH = 16}, và độ rộng thanh ghi tích lũy \texttt{ACC\_WIDTH = 40}. Các cổng giao tiếp gồm xung nhịp \texttt{clk}, xóa \texttt{rst\_n}, cờ cấu hình chuyển mạch \texttt{bypass}, kèm theo cặp giao tiếp Avalon-ST Sink (ngõ vào) và Source (ngõ ra).
    \item \textbf{Dòng 43--49:} Định nghĩa 3 trạng thái của FSM: \texttt{S\_IDLE} (chờ mẫu âm thanh mới), \texttt{S\_CALC} (thực hiện 16 bước nhân tích lũy gấp thời gian), và \texttt{S\_WAIT\_OUT} (giữ mẫu đầu ra chờ hạ nguồn tiếp nhận khi có phản áp). Khai báo thanh ghi bước \texttt{step} 4-bit ($0..15$).
    \item \textbf{Dòng 52:} Khai báo mảng 31 thanh ghi dịch \texttt{shift\_reg[0:30]}, lưu giữ 31 mẫu âm thanh gần nhất.
    \item \textbf{Dòng 55--60:} Khai báo thanh ghi tích lũy \texttt{acc} 40-bit có dấu, thanh ghi đệm ngõ ra \texttt{out\_data\_reg} 16-bit và cờ báo hợp lệ \texttt{out\_valid\_reg}.
    \item \textbf{Dòng 62--70:} Định nghĩa hai điều kiện bắt tay \texttt{s\_handshake} và \texttt{m\_handshake}. Biểu thức gán tín hiệu sẵn sàng \texttt{sink\_ready = (!out\_valid\_reg || source\_ready)} khi ở trạng thái \texttt{S\_IDLE} hoặc khi bật \texttt{bypass}, phản ánh chính xác chuẩn giao tiếp \texttt{readyLatency = 0}.
    \item \textbf{Dòng 74--92:} Khởi tạo bảng ROM chứa 16 hệ số đối xứng Q1.15 độc nhất. Hệ số trung tâm $h[15]$ có giá trị lớn nhất là \texttt{16'sh7ccd} ($31949$), các hệ số còn lại dao động từ $504$ đến $818$ với dấu đổi luân phiên.
    \item \textbf{Dòng 96--100:} Bộ tiền cộng Pre-adder tổ hợp 17-bit: lấy mẫu \texttt{tap\_a = shift\_reg[step]} cộng với mẫu đối xứng \texttt{tap\_b = shift\_reg[30 - step]}. Khi \texttt{step == 15}, \texttt{tap\_b} được ép về $0$ để bảo toàn giá trị mẫu trung tâm $x[n-15]$.
    \item \textbf{Dòng 101--106:} Bộ nhân tổ hợp $17\text{ bit} \times 16\text{ bit} \to 33\text{ bit}$ sinh tích \texttt{prod}, sau đó mở rộng dấu 7 bit thành 40-bit (\texttt{prod\_ext}). Đường dây \texttt{acc\_final} tính trước giá trị tích lũy mới trong chu kỳ xung nhịp hiện tại.
    \item \textbf{Dòng 108--113:} Mạch tổ hợp làm tròn bán phần (\texttt{+ 40'sd16384}) kết hợp dịch phải số học 15 bit và khâu bão hòa chống tràn về khoảng $[-32768, 32767]$.
    \item \textbf{Dòng 117--126:} Nhánh reset bất đồng bộ tích cực mức thấp: xóa toàn bộ 31 thanh ghi dịch, thanh ghi tích lũy và các cờ điều khiển về trạng thái 0 an toàn.
    \item \textbf{Dòng 127--144:} Nhánh xử lý chế độ \texttt{bypass}: mạch hoạt động như một tầng đệm đường ống (pipeline stage) 1 chu kỳ. Khi có giao dịch ngõ vào, dữ liệu \texttt{sink\_data} được chốt thẳng sang \texttt{out\_data\_reg} mà không qua tính toán FIR, thông lượng đạt tối đa 1 từ/chu kỳ.
    \item \textbf{Dòng 145--192:} Máy trạng thái xử lý lọc FIR gấp thời gian:
    \begin{itemize}
        \item Tại \texttt{S\_IDLE}: Khi có mẫu mới (\texttt{s\_handshake} = 1), mẫu mới được nạp vào đầu mảng dịch, toàn bộ 30 mẫu cũ dịch sang phải 1 vị trí. FSM nhảy sang \texttt{S\_CALC}, khởi tạo \texttt{step = 0} và xóa \texttt{acc} về 0.
        \item Tại \texttt{S\_CALC}: Bộ đếm \texttt{step} tăng từ 0 đến 15. Bộ tích lũy \texttt{acc} cộng dồn từng tích số. Đến bước 15, giá trị \texttt{sat\_out} được chốt vào thanh ghi ngõ ra \texttt{out\_data\_reg} và cờ \texttt{out\_valid\_reg} bật lên 1. Nếu hạ nguồn đang sẵn sàng (\texttt{source\_ready} = 1), FSM quay ngay về \texttt{S\_IDLE}; nếu hạ nguồn đang bận, FSM chuyển sang \texttt{S\_WAIT\_OUT}.
        \item Tại \texttt{S\_WAIT\_OUT}: Mạch đóng băng ngõ vào (\texttt{sink\_ready} = 0) và giữ nguyên mẫu đầu ra cho đến khi hạ nguồn tiếp nhận thành công.
    \end{itemize}
\end{itemize}

\subsubsection*{Bảng đối chiếu song sinh: Mã Python Twin so với mã phần cứng Verilog}
Để khẳng định tính tương đương toán học và chu kỳ tuyệt đối giữa mô hình phần mềm và phần cứng, Bảng~\ref{tab:ch03_twin_vs_hdl} đối chiếu chi tiết từng khối lệnh trong tệp mô hình song sinh \texttt{verify\_v3/ch03\_fir\_twin.py} với mã nguồn phần cứng tương ứng:

\begin{table}[htbp]
\centering
\small
\setlength{\tabcolsep}{4pt}
\caption{Đối chiếu trực tiếp giữa các dòng lệnh của Python Twin và mã nguồn phần cứng Verilog}
\label{tab:ch03_twin_vs_hdl}
\begin{tabularx}{\textwidth}{ccX}
\toprule
Dòng Twin (\texttt{ch03\_fir\_twin.py}) & Dòng HDL (\texttt{ch03\_fir\_filter.v}) & Ý nghĩa kỹ thuật tương đương \\
\midrule
Dòng 30--38 & Dòng 117--126 & Thiết lập trạng thái ban đầu khi reset tích cực (\texttt{rst\_n} = 0) \\
Dòng 41--47 & Dòng 66--67 & Logic tổ hợp xác định cờ sẵn sàng \texttt{sink\_ready} theo chuẩn \texttt{readyLatency}=0 \\
Dòng 50--51 & Dòng 62--63 & Đánh giá điều kiện hoàn tất bắt tay \texttt{s\_handshake} và \texttt{m\_handshake} \\
Dòng 54--66 & Dòng 127--143 & Xử lý chế độ chuyển mạch \texttt{bypass} với độ trễ chốt 1 chu kỳ \\
Dòng 68--75 & Dòng 147--161 & Trạng thái \texttt{S\_IDLE}: chốt mẫu vào mảng dịch và chuyển sang \texttt{S\_CALC} \\
Dòng 80--82 & Dòng 96--100 & Rút cặp mẫu đối xứng và thực hiện phép tiền cộng 17-bit \\
Dòng 83--85 & Dòng 101--106 & Nhân hệ số Q1.15 sinh tích 33-bit và cộng dồn vào bộ tích lũy 40-bit \\
Dòng 87--95 & Dòng 168--178 & Bước 15: hoàn tất tính toán, làm tròn $+16384$, bão hòa và bật cờ ngõ ra \\
Dòng 96--98 & Dòng 183--188 & Trạng thái \texttt{S\_WAIT\_OUT}: giữ ngõ ra khi gặp phản áp hạ nguồn \\
\bottomrule
\end{tabularx}
\end{table}

\subsubsection*{Vì sao khối HDL khớp hoàn toàn với khối IP của Intel?}
Khi đặt bản thiết kế HDL tự viết cạnh khối IP FIR Compiler II được cấu hình ở mục 3.3, hai thiết kế khớp nhau hoàn toàn về mặt số học và giao tiếp:
\begin{enumerate}[leftmargin=1.5em, itemsep=3pt]
    \item \textbf{Cơ chế tiền cộng Pre-adder:} Biểu thức \texttt{pre\_add = \$signed(tap\_a) + \$signed(tap\_b)} ở dòng 99 khớp chính xác với phần cứng Pre-adder mà tùy chọn \texttt{Symmetry = Symmetrical} kích hoạt trong khối IP. Cả hai đều triệt tiêu được 15 phép nhân số học.
    \item \textbf{Cơ chế làm tròn bán phần:} Phép toán \texttt{(acc\_final + 40'sd16384) >>> 15} ở dòng 109 tái hiện chính xác chế độ \texttt{Output LSB Rounding = Nearest/Convergent} của khối IP, bảo toàn tính đối xứng của đáp ứng xung và không tạo ra điện áp trôi DC.
    \item \textbf{Khâu bão hòa biên độ:} Mạch kiểm tra vượt ngưỡng $[-32768, 32767]$ ở dòng 110--112 khớp hoàn toàn với tùy chọn \texttt{Output MSB Rounding = Saturation} của IP, ngăn chặn tuyệt đối hiện tượng cuộn bit âm khi tín hiệu đạt đỉnh.
    \item \textbf{Giao thức bắt tay Avalon-ST:} Khối HDL hỗ trợ đầy đủ cơ chế phản áp tổ hợp (tín hiệu \texttt{sink\_ready} sẵn sàng khi ngõ ra rỗng hoặc hạ nguồn bật \texttt{source\_ready = 1}) ở chuẩn \texttt{readyLatency = 0}, tương thích 100\% với giao thức streaming chuẩn của Intel.
\end{enumerate}

\subsubsection*{Đổi lựa chọn thì sao: Đánh đổi giữa các kiến trúc FIR khác nhau}
Nếu thay đổi yêu cầu thiết kế sang một bài toán khác, việc điều chỉnh kiến trúc HDL sẽ tác động ra sao đến các chỉ số hiệu năng?
\begin{itemize}[leftmargin=1.5em, itemsep=3pt]
    \item \textbf{Chuyển sang kiến trúc song song trực tiếp (Fully Parallel FIR):} Nếu bài toán đòi hỏi xử lý dòng âm thanh tốc độ cao (ví dụ xử lý đa kênh hoặc siêu âm 50 MSPS), ta có thể bố trí 16 bộ nhân DSP và một cây cộng 4 tầng thanh ghi. Khi đó, thông lượng tăng vọt lên 1 mẫu/chu kỳ (tăng 17 lần), độ trễ đường ống cố định khoảng 4 chu kỳ. Đổi lại, tài nguyên DSP tiêu tốn tăng từ 1 khối lên 16 khối DSP (chiếm $14{,}3\%$ tổng số DSP của chip Cyclone V).
    \item \textbf{Chuyển sang kiến trúc chuyển vị (Transposed Form):} Cấu trúc chuyển vị đặt các thanh ghi trễ nằm xen kẽ giữa các bộ cộng ngõ ra. Ưu điểm của kiến trúc này là không cần cây cộng lớn ở ngõ ra, giúp đẩy tần số cực đại $F_{\max}$ lên rất cao ($> 200\text{ MHz}$). Tuy nhiên, cấu trúc chuyển vị không thể tận dụng được bộ tiền cộng Pre-adder đối xứng, do đó bắt buộc phải sử dụng đủ 31 bộ nhân, làm tăng gấp đôi chi phí phần cứng.
\end{itemize}

Đối với bài toán ước lượng $F_0$ với tần số lấy mẫu 16 kHz, kiến trúc gấp thời gian 1 DSP với bộ tiền cộng đối xứng mà chúng tôi hiện thực hóa là giải pháp tối ưu nhất xét trên mọi phương diện: diện tích silicon, độ trễ và thông lượng.

\textit{Tóm lại: Khối HDL folded FIR dùng đúng 1 DSP và bộ tiền cộng 17 bit để tính 31 tap trong 16 chu kỳ; bộ tích lũy 40 bit kết hợp làm tròn bán phần (+16384) và bão hòa đảm bảo khớp từng bit với cấu hình IP chuẩn.}



